% ============================================================================
% Chapter 8: Flutter Integration and Headless Execution
% ============================================================================
\section{Flutter Integration and Headless Execution}
\label{sec:flutter-integration}

Traditional web servers respond to HTTP requests over the network. This model is proven and utilized for building cloud architectures (see Chapter \ref{sec:architecture}). However, what happens when we desire to translate the same business logic and computational power into an embedded environment such as a cross-platform mobile or desktop application?

Coroute solves this problem through a fully integrated ``headless'' (without network interface) mode, which interacts with a Flutter frontend via a Foreign Function Interface (FFI). This approach allows the Flutter UI to communicate directly with the C++ core within the same process, bypassing the network stack entirely.

\subsection{Interaction between C++ and Dart (FFI)}
High-level integration requires the creation of a common communication layer between the system language (C++) and the user interface (Dart). For this purpose, the \texttt{bridge/} module was created, exporting a clean C API:

\begin{lstlisting}[language=C++, basicstyle=\small\ttfamily, caption={Exporting C functions from Coroute}]
extern "C" {
    COROUTE_EXPORT void* coroute_init(const char* config_json);
    COROUTE_EXPORT void coroute_shutdown(void* engine_ptr);
    COROUTE_EXPORT void coroute_dispatch_request(
        void* engine_ptr, 
        const char* method, 
        const char* path, 
        const char* body, 
        void (*callback)(int, const char*)
    );
}
\end{lstlisting}

In this architecture, the Flutter application initializes an \texttt{App} object in C++ memory and uses the returned pointer as context. The \texttt{coroute\_dispatch\_request} method allows Flutter to simulate HTTP requests internally in memory. This means all existing C++ endpoints execute unchanged, and the response is asynchronously returned to Dart via the provided callback pointer.

\subsection{Integration Tooling}
Integrating a large CMake project like Coroute inside a Flutter project structure can be complex and cumbersome due to the specifics of the build systems for Android, iOS, Windows, Linux, and macOS.

To automate this integration, Coroute provides specialized CMake configurations: \texttt{coroute\_framework} and \texttt{CorouteApp.cmake}.

\subsubsection{\texttt{coroute\_framework}}
This CMake target gathers all the functionality of the web framework and packages it as a dynamic (Shared) Library --- \texttt{libcoroute.so}, \texttt{coroute.dll}, or \texttt{libcoroute.dylib}, depending on the operating system.

\subsubsection{\texttt{CorouteApp.cmake} and the \texttt{.flutter} Sandbox}
This auto-generated sandbox manages the initialization of the Flutter project. Upon calling \texttt{add\_coroute\_app()}, CMake:
\begin{itemize}
    \item Outlines the required source files for the C++ business logic.
    \item Automatically links \texttt{coroute\_framework}.
    \item Writes a \texttt{.coroute\_lib\_path} configuration file consumed by the Dart build hook to locate the compiled library and the Flutter project root.
    \item Patches the macOS \texttt{Runner/Release.entitlements} and \texttt{Runner/DebugProfile.entitlements} to unconditionally enable \texttt{com.apple.security.network.client}, which \texttt{flutter create} omits by default.
\end{itemize}

\subsubsection{Dart Build Hook (\texttt{hook/build.dart})}
The \texttt{coroute\_framework} package provides a Dart build hook compliant with the Dart 3.10+ Native Assets protocol (\texttt{package:hooks}, \texttt{package:code\_assets}). The hook runs automatically as part of every \texttt{flutter run} or \texttt{flutter build} invocation and performs two functions:

\begin{enumerate}
    \item \textbf{Library bundling}: scans all \texttt{.coroute\_lib\_path.<TARGET\_NAME>} config files in the package root, identifies the one belonging to the current Flutter project by deriving the project root from \texttt{input.outputDirectory} (six levels up through \texttt{.dart\_tool/hooks\_runner/}), copies the pre-built \texttt{libcoroute\_app} into the hook output directory, and declares a \texttt{CodeAsset(DynamicLoadingBundled())}. The Flutter toolchain then bundles and code-signs the library automatically on all five platforms. Using per-project config files (keyed by \texttt{TARGET\_NAME}) ensures that concurrent builds of multiple Coroute projects do not overwrite each other's configuration.
    \item \textbf{Permission injection}: patches \texttt{android/app/src/main/AndroidManifest.xml} (INTERNET permission) idempotently. macOS entitlements are patched at CMake configure time via \texttt{file(WRITE \ldots)} rather than inside the Xcode build phase, as Xcode prohibits modification of entitlements files during a build.
\end{enumerate}

\subsubsection{Static OpenSSL Linking on Apple Platforms}
Because \texttt{libcoroute\_app.dylib} links OpenSSL for TLS support, and macOS App Sandbox blocks access to absolute Homebrew paths (e.g., \texttt{/opt/homebrew/opt/openssl@3/lib/}) at runtime, \texttt{CorouteApp.cmake} explicitly links the OpenSSL static archives (\texttt{libssl.a}, \texttt{libcrypto.a}) directly into the shared library on all Apple targets. The \texttt{-Wl,-dead\_strip\_dylibs} flag removes the residual dynamic stubs that would otherwise arrive transitively. The result is a self-contained \texttt{libcoroute\_app.dylib} (and its \texttt{.framework} wrapper) with no external dependencies blocked by the sandbox.

This approach removes the necessity for the programmer to handle complex FFI configuration build properties at the OS level, providing a true ``plug-and-play'' experience when developing cross-platform applications featuring C++ backend logic and Flutter UI.

\subsection{Event Management across the Language Boundary}
One of the most significant challenges in FFI is managing asynchronous events. As defined in Chapter \ref{sec:architecture}, Coroute relies on C++20 coroutines (\texttt{Task<T>}), which are executed in a dedicated \texttt{IoContext} thread pool.

Concurrently, the Flutter UI operates in its own main isolate (thread). To avoid data races and UI thread blocking, every invocation from Dart to C++ is asynchronous:
\begin{enumerate}
    \item Dart calls \texttt{coroute\_dispatch\_request} by creating a C-struct.
    \item The C++ code launches a coroutine and immediately yields control back to the Dart isolate.
    \item The coroutine executes in the background within the \texttt{IoContext}.
    \item When the coroutine finishes and produces a result, it invokes the provided C-callback.
    \item The FFI layer in Dart translates this result back into a Dart \texttt{Future<String>} or equivalent object that Flutter can visualize.
\end{enumerate}

This coexistence of two distinct asynchronous models (C++ coroutines and Dart Futures) across a tight C FFI boundary is the foundation of the library's capability to build high-performance mobile applications.
